\begin{lemma}[Comparison of the two quadratic forms]
Let $E$ be a finite-dimensional real inner-product space and let
$P,P^{-1}:E\to E$ be continuous linear maps with $P^{-1}(Pz)=z$ for all $z\in E$.
Write $Q=\mathrm{id}-P$ and $A_0=QP^{-1}$. Let $c>0$ satisfy
$c|z|\le|Pz|$ for all $z\in E$, let $\delta_1$ satisfy
$0<\delta_1\le c/2$, and let $h,k\in E$ satisfy $|k-Ph|\le\delta_1|h|$. Then
\[
  |h|\le\frac{2}{c}|k|
  \qquad\text{and}\qquad
  \bigl|\langle Qh,Ph\rangle-\langle A_0k,k\rangle\bigr|
    \le C_1\delta_1|k|^2,
\]
where
\[
  C_1=\frac{4\,\|Q\|\,\|P^{-1}\|\,\|P\|}{c^2}+\frac{2\,\|A_0\|}{c}.
\]
\end{lemma}
\begin{proof}
\emph{The bound on $|h|$.}
By the triangle inequality, the lower bound on $P$ and the hypothesis on $k$,
\[
  |k|\ge|Ph|-|k-Ph|\ge c|h|-\delta_1|h|=(c-\delta_1)|h|\ge\frac{c}{2}|h|,
\]
the last step because $\delta_1\le c/2$. Since $c>0$ this gives
\begin{equation}\label{eq:pqfc-h}
  |h|\le\frac{2}{c}|k| .
\end{equation}

\emph{Two consequences.}
Combining \eqref{eq:pqfc-h} with the hypothesis $|k-Ph|\le\delta_1|h|$,
\begin{equation}\label{eq:pqfc-one}
  |k-Ph|\le\delta_1|h|\le\frac{2\delta_1}{c}|k| .
\end{equation}
Applying $P^{-1}$ to $Ph-k$ and using $P^{-1}(Ph)=h$ together with the operator
norm bound,
\begin{equation}\label{eq:pqfc-two}
  \bigl|h-P^{-1}k\bigr|=\bigl|P^{-1}(Ph)-P^{-1}k\bigr|
   =\bigl|P^{-1}(Ph-k)\bigr|\le\|P^{-1}\|\,|Ph-k|
   \le\|P^{-1}\|\frac{2\delta_1}{c}|k| ,
\end{equation}
using \eqref{eq:pqfc-one} and $|Ph-k|=|k-Ph|$ at the last step.

\emph{The algebraic identity.}
Inserting and removing the term $\langle QP^{-1}k,Ph\rangle$, and using
$A_0=QP^{-1}$ and the bilinearity of the inner product,
\[
\begin{aligned}
  \langle Qh,Ph\rangle-\langle A_0k,k\rangle
   &=\bigl[\langle Qh,Ph\rangle-\langle QP^{-1}k,Ph\rangle\bigr]
     +\bigl[\langle QP^{-1}k,Ph\rangle-\langle QP^{-1}k,k\rangle\bigr]\\
   &=\bigl\langle Q\bigl(h-P^{-1}k\bigr),Ph\bigr\rangle
     +\bigl\langle QP^{-1}k,\,Ph-k\bigr\rangle ,
\end{aligned}
\]
where the first bracket uses the linearity of $Q$ and the second the linearity of
the inner product in its second argument.

\emph{Estimating the two terms.}
From \eqref{eq:pqfc-h} and the operator norm bound,
$|Ph|\le\|P\|\,|h|\le(2\|P\|/c)|k|$. Hence by the Cauchy--Schwarz inequality,
the operator norm bound for $Q$ and \eqref{eq:pqfc-two},
\[
  \bigl|\bigl\langle Q\bigl(h-P^{-1}k\bigr),Ph\bigr\rangle\bigr|
   \le\|Q\|\,\bigl|h-P^{-1}k\bigr|\cdot|Ph|
   \le\|Q\|\cdot\|P^{-1}\|\frac{2\delta_1}{c}|k|\cdot\frac{2\|P\|}{c}|k|
   =\frac{4\,\|Q\|\,\|P^{-1}\|\,\|P\|}{c^2}\,\delta_1|k|^2 .
\]
Likewise, by Cauchy--Schwarz, $|A_0k|\le\|A_0\|\,|k|$ and \eqref{eq:pqfc-one},
\[
  \bigl|\bigl\langle QP^{-1}k,\,Ph-k\bigr\rangle\bigr|
   =\bigl|\langle A_0k,\,Ph-k\rangle\bigr|
   \le\|A_0\|\,|k|\cdot\frac{2\delta_1}{c}|k|
   =\frac{2\,\|A_0\|}{c}\,\delta_1|k|^2 .
\]
Adding the two estimates and using the triangle inequality on the identity above
gives
\[
  \bigl|\langle Qh,Ph\rangle-\langle A_0k,k\rangle\bigr|
   \le\Bigl(\frac{4\,\|Q\|\,\|P^{-1}\|\,\|P\|}{c^2}+\frac{2\,\|A_0\|}{c}\Bigr)
      \delta_1|k|^2=C_1\delta_1|k|^2 ,
\]
which together with \eqref{eq:pqfc-h} is the assertion.
\end{proof}
